% 10_qemu_npu
\section{QEMU RKNPU 功能级仿真与模型结构协同优化}

RKNPU 是 RK3588 上的神经网络加速器，其寄存器接口、命令编码与算子语义均无公开文档，RKNN 运行时以闭源二进制驱动它。此前对这块加速器的驱动 bring-up、RKNN 的持续集成与模型结构调优只能在开发板上进行，每一次改动都要占用一台真机，逐算子的成本也无从读出。功能级 QEMU 设备模型在模拟器中重建这块加速器的寄存器、命令流水线与算子执行，在无板环境下算出与真机一致的计算输出，把闭源加速器变成一个寄存器精确、可进 CI、可逐算子度量的开发与优化载体。此非一个仅接受写入而不计算的桩：模拟器真实执行 INT8 与浮点卷积并回写正确张量。仿真侧现已\ours{实现完整的功能级 RK3588 RKNPU 模拟器}，以两个合并请求（qemu\#19、qemu\#26）合入 \texttt{processmission/qemu}，配套 106 个 qtest 回归，接入 Rockchip IOMMU，并支持三核协同会合；网球检测模型的结构协同优化在该模拟器上完成。本节先说明 RKNPU 与功能级仿真的必要性，再逐块说明命令流水线、IOMMU 接入、算子执行、回归验证与多核协同的机制与实现，最后给出基于该模拟器的模型结构协同优化。本节的每一个模型数字均为该模拟器的逐推理实测，属结构成本与运行时计数口径，非真机读数。

\figphl{../figures/out/Fqemu.png}{功能级 RKNPU 模拟器的分层结构，青色为仿真侧新增的设备模型}{fig:qemu-npu-arch}

图~\ref{fig:qemu-npu-arch} 给出该模拟器的分层结构。既有的部分是 QEMU 主线提供的设备模型框架、标准 IOMMU 的 \texttt{MemoryRegion} 翻译接口、迁移状态机制与 qtest 基础设施。新增的部分是 RK3588 RKNPU 的设备接入（MMIO、中断、复位、设备树绑定）、命令流水线各功能块的寄存器与状态机、真实的算子执行路径、多核会合逻辑，以及为模型优化服务的逐推理 profiling 计数。RKNN 运行时二进制在模拟器上未经改动直接运行，模拟器对齐其可观察的寄存器与命令行为到足以让整条推理链正确执行。

\subsection{RKNPU 与功能级仿真的必要性}

RKNPU 由一条命令驱动的定点/浮点计算流水线构成，RKNN 运行时把编译好的模型翻译成一串寄存器写入与命令表，提交给加速器执行卷积、激活、池化等算子并回收输出张量。由于寄存器布局与命令编码闭源，任何脱离开发板的开发都无从下手：驱动无法在 CI 中回归，模型改动无法在无板环境下验证，逐算子的成本更无法读出。功能级仿真的价值在于同时满足两点。其一是寄存器精确，模拟器复现命令流水线的可观察行为，使未经改动的 RKNN 运行时能在其上运行整条推理链。其二是功能正确，模拟器在维护寄存器状态之外据命令算出真实的 INT8 与浮点输出，因而可用作精度评估的参考。二者合起来使这块加速器可脱离开发板被调试、被度量、被优化。

\subsection{设备接入与命令流水线}

模拟器以 \texttt{-machine rk3588-evb,rknpu=on} 启用，机器随之发出配套的设备树绑定，向访客暴露 RKNPU 的 MMIO 区、中断线与复位。命令处理沿加速器的功能块展开：RKNN 运行时把张量与权重的地址、算子参数与协同模式写入各块的寄存器，再以命令表把任务串成任务链提交，模拟器按块推进状态机，从核提交、寄存器待决状态、复位与迁移状态一并建模。各功能块的职责如下。

\paragraph{PC 程序控制器}
PC 是命令流水线的入口，负责从命令表取出任务、解出各算子的寄存器配置并驱动其余各块，同时维护任务链与待决寄存器状态。模拟器重建这条取指与派发路径，使一次推理对应的整串任务按序推进。

\paragraph{CNA 卷积前端}
CNA 负责卷积的输入组织，把特征图与权重按卷积核的取数模式喂入计算核，并承载协同模式的配置位。模拟器复现其取数几何与配置寄存器，其中协同模式位（见多核协同一节）即由 CNA 的控制寄存器解出。

\paragraph{CORE 卷积计算核}
CORE 是乘累加阵列，对喂入的特征与权重做卷积与深度卷积的乘累加。模拟器在此真实执行定点与浮点的乘累加，覆盖 INT8、量化 INT8 与 FP16/FP32，算出的部分和交由后处理。

\paragraph{DPU 后处理}
DPU 承接卷积输出，完成激活、量化、逐元素运算与查表（LUT）等后处理。激活函数中形如 SiLU 的非线性会被编译为查表算子在此执行，模拟器复现 LUT 与逐元素路径的功能行为，这一点直接关系到下文模型优化中激活函数替换的成本变化。

\paragraph{DPU-RDMA 读通道}
DPU-RDMA 是后处理块的读 DMA，将待处理张量从内存读入 DPU。模拟器建模其读取几何与面选择，多核路径下的双端口面选择与有效面缺口布局在此处理。

\paragraph{PPU 池化}
PPU 负责池化与平面处理。模拟器复现池化算子及其输出布局，与卷积、后处理一同构成完整的算子集合。

\subsection{IOMMU 接入}

RKNPU 经 Rockchip IOMMU 访问访客内存，模拟器通过 QEMU 标准 IOMMU 的 \texttt{MemoryRegion} 接口接入这条翻译路径。加速器发出的每一次 DMA 都经 IOMMU 做地址翻译与权限检查，翻译失败或权限不足时如实上报 DMA 翻译错误，不会静默访问宿主内存。此接入使模拟器可验证驱动侧的地址映射是否正确，也把畸形地址的越界访问挡在翻译层。

\subsection{算子的功能执行}

功能执行路径是模拟器中难度最高的一环，模拟器在维护寄存器状态之外\ours{据命令算出真实输出}。执行覆盖 INT8、量化 INT8 与 FP16/FP32 三种数值路径，算子覆盖卷积、深度卷积、逐元素、查表、池化以及张量布局转换。量化 INT8 路径按加速器的定点约定做乘累加与再量化，浮点路径按 FP16/FP32 计算，查表路径复现激活函数编译所得的 LUT 语义，布局转换复现加速器在各块之间对张量摆放格式的重排。因输出与真机一致，该路径既支撑驱动 bring-up 的正确性验证，也可直接作为模型精度评估的参考基准。为拦截畸形或超大的访客负载，模拟器对宿主内存占用与单次算子规模设定了可配上限，越限即被安全地拒绝而不致耗尽宿主资源。

\subsection{qtest 回归}

模拟器附带 106 个 qtest 子测试，即机器可判定的判定门，一次运行即回归整套设备行为与算子正确性。覆盖面包括 MMIO、复位、IRQ 与 IOMMU 行为，命令解码与任务生命周期，整数与浮点执行的数值正确性，DMA 错误与不支持控制字的处理，任务链，迁移状态，以及 DPU-RDMA 与 PPU 的布局。这套回归使模拟器的每一处行为改动都可自动核验，是把它用作 RKNN 的 CI 与精度参考的前提。上述设备接入、IOMMU、算子执行与 qtest 以合并请求 qemu\#19 合入 \texttt{processmission/qemu}，已合并。

\subsection{多核协同}

RK3588 的 RKNPU 由三个物理核构成，卷积任务可按协同模式在一至三核上并行并在完成边界会合。模拟器对照真实硬件\ours{逆向出这套协同行为并加以复现}。协同模式自 CNA 的 \texttt{CONV\_CON2} 寄存器 [30:28] 位解出：模式 0 为单核，模式 1 与 4 为双核，模式 2 与 5 为三核，模式 3、6、7 另有区别。每个 NPU 核在模拟器中运行于独立的宿主 worker 上，各自推进本核的任务，在协同完成边界处同步。逆向所依据的硬件证据一并记录：双核模式 4 等待 core0 与 core1，三核模式 2 与 5 等待 core0 至 core2；参与核数相同的模式即便模式值不同也在同一点会合；各核本地命令表的任务下标无须跨核对齐；只要物理核集合与参与核数属同一类，独立分配的命令链与输出缓冲亦满足同一硬件会合条件。扩展至多核时一并修正了若干数值正确性问题：INT16 unpool 的几何与存储映射、FP16 DPU-RDMA 的双端口面选择与有效面缺口布局、通道分组下 FP16 深度卷积权重的紧凑索引，以及跨步输出回写仅落在有效输出通道面上。这部分以合并请求 qemu\#26 合入 \texttt{processmission/qemu}，已合并。下文模型优化所依赖的逐推理 profiling 计数由一个在审的合并请求提供。

\subsection{基于模拟器的模型结构协同优化}

模拟器暴露的逐算子成本计数，使网球检测器得以逐层修改结构，每一步改动都可读出 MAC、命令字、任务取指、LUT 任务数与逐段耗时的变化，这是板级黑盒跑分无法提供的信息。度量口径固定为 RKNN 运行时 2.3.2，精度在 285 图 / 382 框的保留集上按 INT8 评估，置信度 0.25，NMS IoU 0.45。表~\ref{tab:npu-ladder} 给出累计的结构成本，每一行均在上一行基础上再叠加一层改动。

\begin{table}[H]\centering\small
\resizebox{\linewidth}{!}{%
\begin{tabular}{@{}lrrrrrr@{}}
\toprule
阶梯步骤 & MAC (GMac) & RDMA 输入 (M) & 任务取指 & 命令字 (k) & LUT 任务 & 逻辑输出 (KiB) \\
\midrule
原始 640$\times$640 SiLU & 3.15 & 27.21 & 443 & 95.15 & 107 & 541 \\
输入缩放 480$\times$640 & 2.36 & 20.41 & 422 & 92.42 & 86 & 406 \\
改用 ReLU & 2.36 & 19.60 & 502 & 43.56 & \textbf{3} & 406 \\
去掉 score-sum（9$\to$6 输出） & 2.36 & 19.50 & 484 & 42.62 & 3 & 400 \\
类别头 64$\to$32 & 2.07 & 19.09 & 484 & 42.62 & 3 & 400 \\
头框合并（隐藏 32 / reg\_max 8） & \textbf{2.06} & \textbf{18.89} & 484 & \textbf{42.62} & 3 & \textbf{203} \\
\bottomrule
\end{tabular}}
\caption{模拟器实测的累计结构成本，逐步压缩 MAC、命令字与逻辑输出}
\label{tab:npu-ladder}
\end{table}

对全链路成本影响最大的一步是\ours{激活函数的替换}。原始 640$\times$640 模型采用 SiLU，SiLU 在 DPU 上被编译为大量查表算子。输入缩放后单次推理含 86 个 LUT 任务，\texttt{rknn\_run} 为 12.31 ms；替换为 ReLU 后 LUT 任务降至 3 个，命令字由 92.42k 降到 43.56k，\keyres{\texttt{rknn\_run} 由 12.31 ms 减半至 5.79 ms}。其余各步分别修正一处结构冗余。输入由 640$\times$640 调整为模型原生的 480$\times$640，节省 RDMA 输入与 MAC。去掉冗余的 score-sum 将 9 路输出并为 6 路。类别头由 64 收窄至 32。检测头与框回归头合并，隐藏维取 32、reg\_max 取 8，逻辑输出随之由 400 收窄至 203 KiB。上述改动累计将 MAC 由 3.15 压缩至 2.06 GMac，逻辑输出由 541 收窄至 203 KiB，RDMA 输入与命令字一并下降。

精度在整条阶梯上保持稳定。如表~\ref{tab:npu-accuracy}，mAP50 由 92.9\% 变化至 91.1\%，中途最低降至 89.5\%，mAP50-95 略有上升，由 69.9\% 至 70.6\%；Precision 由 93.8\% 至 94.8\%，Recall 由 91.6\% 至 90.3\%。耗时方面 \texttt{rknn\_run} 由 14.66 ms 降到 5.68 ms，\texttt{outputs\_get} 由 0.286 ms 降到 0.125 ms，postprocess 始终保持在 0.03 ms 以下。以可忽略的精度代价把单次推理的运行时耗时\keyres{压缩到约三分之一}（\texttt{rknn\_run} 14.66$\to$5.68 ms），加速器 MAC 降至约三分之二（3.15$\to$2.06 GMac）。

\begin{table}[H]\centering\small
\begin{tabular}{@{}lrrrr@{}}
\toprule
阶梯步骤 & mAP50 (\%) & mAP50-95 (\%) & \texttt{rknn\_run} (ms) & \texttt{outputs\_get} (ms) \\
\midrule
原始 640$\times$640 SiLU & 92.9 & 69.9 & 14.655 & 0.286 \\
输入缩放 480$\times$640 & 91.0 & 69.7 & 12.314 & 0.251 \\
改用 ReLU & 89.5 & 70.0 & \textbf{5.786} & 0.250 \\
去掉 score-sum & 89.5 & 70.0 & 5.839 & 0.157 \\
类别头 64$\to$32 & 90.4 & \textbf{71.6} & 5.643 & 0.154 \\
头框合并 & \textbf{91.1} & 70.6 & 5.677 & \textbf{0.125} \\
\bottomrule
\end{tabular}
\caption{同一阶梯上的 INT8 精度与逐推理耗时，精度基本持平而 \texttt{rknn\_run} 降至约三分之一}
\label{tab:npu-accuracy}
\end{table}

\figph{../figures/out/F13.png}{模拟器实测：网球模型结构协同优化把 MAC 3.15$\to$2.06 GMac、LUT 任务 107$\to$3、\texttt{rknn\_run} 14.66$\to$5.68 ms，同时 mAP50 保持在 92.9$\to$91.1\%}

本节的 MAC、LUT 任务数、\texttt{rknn\_run} 与 mAP50 均为 QEMU RKNPU 功能模拟器的逐推理实测，属结构成本与运行时计数口径，非真机测量。优化后模型在开发板上的端到端首帧与逐帧时延由网球应用的板级运行给出，见启动与推理相关各节。模拟器本身是一项可复用的产物，使这套逐算子的成本与精度分析，连同 RKNN 的 CI 与驱动 bring-up，均可脱离开发板进行。

\figph{../figures/out/Fdemo_npu.png}{仿真 NPU 上运行网球 YOLOv8 推理的真实终端输出：模型加载、INT8 反量化张量与逐帧推理，末端给出检测结果 \texttt{sports ball}、置信度 70.8\%（单命令可复现，无需真机）}
